\subsection{下界的有限证书}
补充材料 \path{paper_claims.json} 保存本稿的准确输入和全分数。
令 $D=10^{35}$、$S=4D$。质量根以式\eqref{pre:root}隔离；事件根 $1$ 取准确值，$1944^{1/8}$ 取相邻的 $D^{-1}$ 有理端点。所有 $32768$ 个标签都重建为固定坐标次序的物理盒，再分别求内、外并集。下表记录足够复核本下界的有限规模。
\begin{center}
\begin{tabular}{lr}
\toprule
项目&数量\\\midrule
带标签源块&$32768$\\
正内盒、正外盒&各 $882$\\
去除完全重复后的物理内盒、外盒&各 $644$\\
准确并集的总缓存状态&$210446$\\
共同半径的全部源坐标包含检查&$34054$\\
\bottomrule
\end{tabular}
\end{center}
完整有理区间位于
\[
\begin{gathered}
 L_0\le 2^8\left|\bigcup B_\ell^-\right|
 \le 2^8\left|\bigcup B_\ell^+\right|\le U_0,\\
 L_0=\frac{21888111794266698273227523980123024}{10^{34}},\\
 U_0=\frac{21888111794266698273227523980123027}{10^{34}}.
\end{gathered}
\]
这里两端是由原始全分数直接向外舍入所得的简化有理界；验证程序仍比较完整内外分数。
取正文中的 $\varepsilon,\theta,K$ 后，程序验证
\[
 \theta\left(L_0-\frac{441}{9536743164062500000}\right)
 \left(\frac{2K+1}{2K+5}\right)^8
 >\frac{2188811179426669}{10^{15}}.
\]

\subsection{上界标量的有理判据}
取 $t_0=3837029/524288$，把 $B_8(t_0)$ 按二项式展开写成
$\alpha+\beta\log t_0$，其中 $\alpha,\beta\in\mathbb Q$。
令 $z=(t_0-1)/(t_0+1)$。使用 $N=400$ 时，
\[
 0\le\log t_0-2\sum_{j=0}^{N-1}\frac{z^{2j+1}}{2j+1}
 \le\frac{2z^{2N+1}}{(2N+1)(1-z^2)}.
\]
按 $\beta$ 的符号传播上下端，得到 $B_8(t_0)\le B_+$，其中可取
\[
 B_+=\frac{133365553664751128154150225918012012532222}{10^{40}}.
\]
记 $s=(1+1/t_0)^4$、$U=56627260639/2000000000$。准确计算满足
\[
 d:=U-s^2-B_+>0,\qquad
 d^2-4s^2B_+>\frac{19}{10^{10}}>0.
\]
因为 $B_+>0$，两次平方保序，由此得到
$F_8(t_0)\le(s+\sqrt{B_+})^2<U$。

\subsection{四维密度的完整锥分解}
为说明应用中点值计算的有限化，取频率排序室
$\xi=(a,b,c,d)$，$a\ge b\ge c\ge d\ge0$。定义射线
\[
\begin{aligned}
 r_1&=(1,0,0,0),&r_2&=(1,1,0,0),&r_3&=(1,1,1,0),\\
 r_4&=(1,1,1,1),&r_5&=(2,1,1,0),&r_6&=(3,1,1,1).
\end{aligned}
\]
排序室按 $a\ge b+c+d$、$a\le b+c-d$ 和中间区域分为
\[
 C_1=\operatorname{cone}(r_1,r_2,r_5,r_6),\quad
 C_2=\operatorname{cone}(r_5,r_2,r_3,r_4),\quad
 C_3=\operatorname{cone}(r_5,r_2,r_4,r_6).
\]
三组锥坐标分别为
\[
\begin{gathered}
 (a-b-c-d,\ b-c,\ c-d,\ d),\\
 (a-b,\ b-c,\ -a+b+c-d,\ d),\\
 \left(c-d,\ b-c,\ \frac{-a+b+c+d}{2},\
                  \frac{a-b-c+d}{2}\right).
\end{gathered}
\]
在相应区域它们非负，且准确重构频率向量，证明三锥覆盖及内部不交。
绝对行列式依次为 $1,1,2$。
对 $\psi_4(\xi)=\sum_{v\in W}|v\cdot\xi|$，六条射线的值依次为
$8,8,12,12,16,24$；每条 $v$ 在一个锥的全部射线上同弱号，故
$\psi_4$ 在该锥上线性。
全部 $2^4\,4!=384$ 个带符号置换给出 $1152$ 个锥，交叠只在零测边界。
因此 Fourier 反演逐锥积分精确给
\[
 p_t^{(4)}(x)=\frac1{16\pi^4}
 \sum_C\frac{|\det(r_{C,1},\ldots,r_{C,4})|}
 {\prod_{j=1}^4\bigl(t\psi_4(r_{C,j})-i x\cdot r_{C,j}\bigr)}.
\]
$t>0$ 时分母均不为零。对 $x=A,B$，将全部项通分后再约分，得到补充数据中的 $P_A,P_B$。
可选程序 \path{supplement/verify_green.py} 重新生成全部锥，对实、虚分子作准确全和，验证约分后的有理函数及三个部分分式积分。零投影和重复投影均保留在全和中处理。

\subsection{复现与材料范围}
\path{supplement/verify_paper.py} 仅依赖 Python 标准库，重建并复算本稿唯一的下界构造，核对有限光滑输入的比值，并以有理数验证上述对数尾和平方判据。它不执行搜索，不计算其他首轴参数，不认证某个构造族的全局最优值。
\begin{verbatim}
python3 -B supplement/verify_paper.py
\end{verbatim}
较快的 \texttt{--scalar-only} 模式只验证保存分数与上界标量，不代替完整并集复算。
\path{green_anchors.json} 保存两个四维有理密度和三个准确积分的部分分式。补充说明另外标明原始程序、来源散列与所需软件；主定理的标准库复核不依赖浮点优化器、MATLAB 或区间库。

正文中的主定理、条件势核推论与点值计算具有不同作用：主定理已给出 $C_8$ 的双侧界；条件势核推论需另有全域 $g_0,I_0$ 证书；点值仅提供校验与必要限制。研究过程记录、版本散列和执行回执不作为额外数学假设。
